\documentclass[reprint, amsmath, amssymb, aps, pra, nofootinbib]{revtex4-2}

% --- Основные пакеты ---
\usepackage[T2A]{fontenc}
\usepackage[utf8]{inputenc}
\usepackage[russian,english]{babel}

\usepackage{graphicx}
\graphicspath{{figures/}}
\usepackage{bm}

% hyperref в REVTeX лучше подключать с параметрами или использовать встроенный механизм
\usepackage[unicode, colorlinks=true, linkcolor=blue, citecolor=blue, urlcolor=blue]{hyperref}

\begin{document}

\title{Название работы}

\author{Ефим Сычев}
\email{sychev.ea@phystech.edu}
\affiliation{Московский физико-технический университет, группа Б02-506}

\author{Иван Рудко}
\email{rudko.ia@phystech.edu} 
\affiliation{Московский физико-технический университет, группа Б02-506}

\date{\today}

\begin{abstract}
Здесь пишется аннотация. Кратко: что измеряли, каким методом, что получили в итоге. Никакой воды, только суть.
\end{abstract}

\maketitle

% --- Основной текст ---
\section{Введение}
Тут начинается твой текст. Если нужно сослаться на формулу, используй команду \ref{eq:my_formula}.

\section{Физические основы явления}
Обычные формулы в строке пишутся так: $E = mc^2$. 
А выделенные формулы оформляются окружением \texttt{equation}:
\begin{equation}
    \int_{-\infty}^{+\infty} e^{-x^2} dx = \sqrt{\pi}
    \label{eq:my_formula}
\end{equation}

\section{Экспериментальная установка}
Здесь мы вставляем картинку.

\section{Ход работы}

\section{Заключение}
Выводы по работе. 

% --- Библиография ---
\begin{thebibliography}{99}
\bibitem{devoret1985}
M. H. Devoret, J. M. Martinis, and J. Clarke, \textit{Phys. Rev. Lett.} \textbf{55}, 1908 (1985).
\end{thebibliography}

\end{document}